\chapter{Unduloids and their geometric limits}\label{sec:unduloid-examples}
The Delaunay unduloids interpolate between cylinders and a singular chain of spheres. The underlying graph family is described in \cite{Maz02}, Section 2.1, and the roulette parametrizations in \cite{BBM13}. The arclength formula below gives direct proofs of completeness, embeddedness, and several limits without invoking a classification of arbitrary complete CMC surfaces.

\begin{theorem}[An explicit arclength construction of unduloids]\label{thm:explicit-unduloid-construction}
Fix $H>0$ and $0<a<1/H$, put $b=2/H-a$ and $c=Hab/2$, and define for all $s\in\mathbb R$
\[
 r_a(s)=\sqrt{a^2+(b^2-a^2)\sin^2(Hs/2)},\qquad
 z_a(s)=\int_0^s\left(\frac H2r_a(t)+\frac c{r_a(t)}\right)dt.
\]
Then $F_a(s,\theta)=(r_a(s)\cos\theta,r_a(s)\sin\theta,z_a(s))$ is a smooth rotational CMC immersion of trace curvature $H$, with $s$ arclength and outward radial normal as in Proposition \ref{prop:rotational-arclength-curvature}. Its neck radius is $a$ and bulge radius $b$. These realize the unduloids of \cite{Maz02}, p.~4, after scaling its $H=1$ normalization.
\end{theorem}
\begin{proof}
All expressions are smooth because $r_a\ge a>0$. With $w=r_a^2$, differentiation gives $(w')^2=H^2(b^2-w)(w-a^2)$. Lemma \ref{lem:unduloid-radius-factorization} therefore implies $(r_a')^2+(z_a')^2=1$. Moreover $z_a'>0$, so the tangent angle can be chosen smoothly with $\sin\psi>0$. The flux $r_a z_a'-Hr_a^2/2$ is identically $c$. On the open intervals where $r_a'\ne0$, differentiation gives $\psi'=H-\sin\psi/r_a$. The zeros of $r_a'$ are discrete, and smoothness extends this identity across them. The curvature formula proves the assertion. The sine formula attains its minimum and maximum, giving the radii. Local uniqueness of the smooth system identifies these solutions with the graph family having the same neck data.
\end{proof}

\begin{proposition}[Periodicity, completeness, and embeddedness]\label{prop:unduloid-global-geometry}
The radius in Theorem \ref{thm:explicit-unduloid-construction} has period $P=2\pi/H$, and
\[
 F_a(s+P,\theta)=F_a(s,\theta)+L_a e_3,\qquad
 L_a=\int_0^P\left(\frac H2r_a+\frac c{r_a}\right)ds>0.
\]
The map on $\mathbb R\times S^1$ is a complete proper embedding. In particular the surface is a global positive cylindrical graph.
\end{proposition}
\begin{proof}
The sine-squared formula gives the radius period and hence the periodicity of $z_a'$, proving the translation formula. The arithmetic-geometric mean inequality and the arclength identity give
$\sqrt{2Hc}\le z_a'\le1$. Thus $z_a$ is strictly increasing and tends to $\pm\infty$ at the respective ends. The height recovers $s$ uniquely, and $r_a>0$ then recovers the angular coordinate, proving injectivity. Bounded ambient sets bound the height and hence $s$; together with compactness of $S^1$, this proves properness. The induced metric $ds^2+r_a^2d\theta^2$ dominates the complete product metric $ds^2+a^2d\theta^2$. An intrinsic Cauchy sequence converges for the product metric, and local smooth equivalence gives convergence in the induced metric. Finally the strictly monotone height gives the claimed global graph.
\end{proof}

\begin{proposition}[The cylinder limit]\label{prop:unduloid-cylinder-limit}
As $a\uparrow1/H$, the parametrizations $F_a$ converge smoothly on compact subsets to the cylinder $(H^{-1}\cos\theta,H^{-1}\sin\theta,s)$. The axial translation lengths satisfy $L_a\to2\pi/H$.
\end{proposition}
\begin{proof}
Then $b\to1/H$, $c\to1/(2H)$, and the explicit radius converges to $1/H$ with all derivatives on compact intervals. The integrand defining $z_a$ converges smoothly to one. Thus $z_a\to s$ smoothly, and integration over the fixed interval $[0,2\pi/H]$ proves the period limit. The limiting cylinder has every axial translation as a symmetry; the limiting value is the limit of the chosen unduloid periods, not a smallest positive period of the cylinder.
\end{proof}

\begin{proposition}[The chain-of-spheres limit]\label{prop:unduloid-sphere-chain-limit}
As $a\downarrow0$, $F_a$ converges uniformly on compact parameter sets to a chain of radius-$2/H$ spheres tangent along the rotation axis. The convergence is smooth away from the neck parameters $s\in(2\pi/H)\mathbb Z$, and $L_a\to4/H$. The limiting image is not a smooth embedded surface at its contact points, and the displayed limiting parametrization is not an immersion there.
\end{proposition}
\begin{proof}
The radius converges uniformly on compact intervals to $(2/H)|\sin(Hs/2)|$, smoothly where the sine is nonzero. On such intervals $c/r_a\to0$ smoothly, so $z_a'\to|\sin(Hs/2)|$. Since $0<z_a'\le1$, dominated convergence gives convergence in $L^1$ on compact intervals and hence uniform convergence of its primitives. On $0\le s\le2\pi/H$ the limit is
\[
 r_0(s)=\frac2H\sin(Hs/2),\qquad z_0(s)=\frac2H(1-\cos(Hs/2)).
\]
Consequently $r_0^2+(z_0-2/H)^2=4/H^2$, and rotation gives one sphere. Its poles are at heights zero and $4/H$. The limiting translation formula repeats it into a tangent chain. Integrating the limiting axial derivative proves $L_a\to4/H$. Two spherical sheets meet at each contact, so the union is not locally an embedded smooth surface; the limiting parametrization collapses an angular circle at each such point. This does not rule out representing the individual spheres by a smooth immersion from a disjoint union of spheres.
\end{proof}

\begin{proposition}[A shrinking neck converges to a catenoid after dilation]\label{prop:unduloid-neck-catenoid-limit}
As $a\downarrow0$, dilate the unduloid by $1/a$ about its neck at height zero and use $s=a\xi$. Then
\[
 a^{-1}F_a(a\xi,\theta)\longrightarrow
 \big(\sqrt{1+\xi^2}\cos\theta,\sqrt{1+\xi^2}\sin\theta,\operatorname{arsinh}\xi\big)
\]
smoothly on compact subsets of $\mathbb R\times S^1$. The limit is the radius-one catenoid in arclength coordinates. The rescaled trace mean curvature is $aH\to0$.
\end{proposition}
\begin{proof}
From the radius formula,
\[
 \frac{r_a(a\xi)^2}{a^2}=1+(b^2-a^2)\left(\frac{\sin(Ha\xi/2)}a\right)^2\longrightarrow1+\xi^2.
\]
The sine quotient extends smoothly to $a=0$ on compact $\xi$ intervals. Taking the positive square root preserves smooth convergence. Write this radius ratio as $R_a$. The derivative in $\xi$ of $z_a(a\xi)/a$ is
\[
 \frac{Ha}{2}R_a+\frac{c/a}{R_a}\longrightarrow\frac1{\sqrt{1+\xi^2}},
\]
since $c/a=Hb/2\to1$. The primitives vanish at zero, so they converge smoothly to $\operatorname{arsinh}\xi$. Substituting $\xi=\sinh v$ gives $(\cosh v\cos\theta,\cosh v\sin\theta,v)$, the catenoid of Example \ref{ex:explicit-catenoid}. Dilation by $1/a$ multiplies principal curvatures by $a$, proving the final assertion.
\end{proof}

\begin{example}[A round sphere is only a local positive cylindrical graph]\label{ex:sphere-zero-flux-graph}
For $H>0$, the graph $r(z)=\sqrt{4/H^2-z^2}$ on $|z|<2/H$ solves the CMC graph equation and has flux zero. It is positive on its domain but is not periodic and does not extend as a positive graph over all $\mathbb R$.
\end{example}
\begin{proof}
Its rotation is the radius-$2/H$ sphere away from its poles, with outward trace curvature $H$. Also $\sqrt{1+r_z^2}=2/(Hr)$, so $r/\sqrt{1+r_z^2}-Hr^2/2=0$. The radius tends to zero at the two finite endpoints. This shows why periodicity statements for positive solutions need their global-domain and flux hypotheses; positivity on an arbitrary interval is insufficient.
\end{proof}

\begin{proposition}[The stability threshold for a periodic cylinder]\label{prop:periodic-cylinder-cmc-index}
In the flat three-manifold $\mathbb R^2\times(\mathbb R/L\mathbb Z)$ with $L>0$, the radius-$R$ cylinder is a closed CMC boundary of trace curvature $1/R$. Its volume-constrained index and nullity are
\[
 \operatorname{ind}_0=2\#\{j\in\mathbb N:j<L/(2\pi R)\},\qquad
 \operatorname{nul}_0=2+2\,\mathbf1_{\{L/(2\pi R)\in\mathbb N\}},
\]
where $\mathbb N=\{1,2,\ldots\}$. Its unrestricted index is $1+\operatorname{ind}_0$. It is volume-constrained stable if and only if $L\le2\pi R$.
\end{proposition}
\begin{proof}
The cylinder bounds the solid torus $B_R^2\times(\mathbb R/L\mathbb Z)$, so the usual volume constraint applies in this ambient quotient. Its metric is $R^2d\theta^2+dz^2$, its principal curvatures are $1/R$ and zero, and the ambient Ricci tensor is zero. Thus the CMC second-variation form is
\[
 Q(f,f)=\int\left(R^{-2}|f_\theta|^2+|f_z|^2-R^{-2}f^2\right)dA.
\]
Fourier expansion in both periodic variables diagonalizes the form. On the complex mode $e^{ik\theta}e^{2\pi ijz/L}$, the eigenvalue of $-J$ is
\[
 \frac{k^2-1}{R^2}+\left(\frac{2\pi j}L\right)^2,\qquad k,j\in\mathbb Z.
\]
The mean-zero condition removes exactly the constant mode. Negative eigenvalues can then occur only when $k=0$ and $0<|j|<L/(2\pi R)$. Each positive $j$ contributes a real sine/cosine pair, proving the constrained index. The always-present zero modes are $\cos\theta,\sin\theta$, induced by horizontal translations. Two further axial modes occur precisely when $j=L/(2\pi R)$ is a positive integer. No other pairs have eigenvalue zero. For the bilinear kernel on the constrained space the same computation applies, since Fourier orthogonality separates the removed constant mode. Reinstating that mode adds one negative direction. The index vanishes precisely at the stated threshold. This is a compact example in the flat quotient, not an embedded compact CMC surface in Euclidean three-space.
\end{proof}

\begin{example}[The infinite cylinder is unstable under compactly supported volume constraints]\label{ex:infinite-cylinder-volume-instability}
The round cylinder $S_R^1\times\mathbb R\subset\mathbb R^3$ has a smooth compactly supported normal speed $f$ with $\int f\,dA=0$ and $Q(f,f)<0$. Thus even the first-order volume constraint does not make the infinite cylinder stable.
\end{example}
\begin{proof}
Choose a nonzero odd smooth function $\eta\in C_c^\infty(\mathbb R)$ and set $f_T(\theta,z)=\eta(z/T)$. The integral of the speed vanishes by oddness. Direct substitution into the cylinder form gives
\[
 Q(f_T,f_T)=2\pi R\left(T^{-1}\int_{\mathbb R}|\eta'|^2
               -\frac T{R^2}\int_{\mathbb R}\eta^2\right)<0
\]
for sufficiently large $T$. The support is compact on the cylinder. If an exactly volume-preserving path is desired, add a second compactly supported normal variation with nonzero integral and solve for its coefficient by the same implicit-function argument as Proposition \ref{prop:realize-volume-constraint}. Here volume means the signed change in the bounded region affected by the compactly supported perturbation, not the infinite total volume of the solid cylinder. The zero first derivative of that correction preserves the displayed negative second variation.
\end{proof}
